
The ability to detect when a large language model has generated a hallucinated response — one that is fluent and confident-sounding but factually incorrect — is a practical problem for safe LLM deployment. Several post-hoc uncertainty quantification (UQ) methods have been developed that derive hallucination detection signals from standard inference outputs, without requiring additional training or labeled data. Three such methods have received significant attention: token-level entropy \citep{Kadavathetal2022}, semantic entropy \citep{Kuhnetal2023}, and SelfCheckGPT-BERTScore \citep{Manakuletal2023}.

Each method has been evaluated separately on different benchmarks. Semantic entropy was validated on TriviaQA and Natural Questions, achieving AUROC ≈ 0.78 \citep{Kuhnetal2023}. SelfCheckGPT was validated on WikiBio biography generation, achieving AUC-PR ≈ 0.85 \citep{Manakuletal2023}. Token entropy has been used as a baseline in multiple studies [Kadavath et al., 2022; Xiao and Wang, 2022]. Because these evaluations use different benchmarks, different models, and different protocols, a practitioner cannot determine from the existing literature which method to use for a specific deployment: short factual QA with an open-source LLM such as LLaMA-2-7B-chat.

This paper addresses that gap by applying all three methods to the same dataset (HaluEval-QA, 2,000 stratified examples), the same model (LLaMA-2-7B-chat), and the same evaluation protocol (AUROC with 1,000-resample bootstrap confidence intervals) under matched inference budgets (N=5 stochastic samples per example). The experimental design controls for confounds present in cross-study comparisons: dataset characteristics, model behavior, and inference variance.

The central result is that none of the three methods achieves above-random hallucination discrimination on HaluEval-QA with LLaMA-2-7B-chat. The failure is not uniform: each method fails for a different identifiable reason. Semantic entropy collapses to a constant signal because the NLI clustering mechanism produces maximum cluster counts (5 out of 5 responses in distinct clusters) for 72.8\% of examples. SelfCheckGPT-BERTScore achieves AUROC = 0.3562, which is significantly below the random baseline of 0.5, indicating a negative — rather than positive — correlation between response consistency and hallucination labels on this benchmark. Token entropy is near-random (AUROC = 0.4829, 95\% CI spanning 0.5).

\textbf{Sub-hypothesis design.} To decompose the mechanism failure, three sub-hypotheses were executed sequentially. H-E1 (EXISTENCE) tested whether a measurable AUROC gap between at least one method pair exists; the gate was satisfied (two pairs with Δ ≥ 0.05 and non-overlapping 95\% CIs), but the gap was driven by SelfCheckGPT performing below random, not by semantic entropy outperforming token entropy. H-M1 (MECHANISM) tested whether the Pearson correlation between token entropy and semantic entropy is below 0.9; semantic entropy was diagnosed as constant (std = 4.14 × 10⁻²⁵), rendering the correlation undefined — the gate was recorded as DEGENERATE_PASS. H-M2 (MECHANISM) tested whether the NLI aggregation rate is at least 0.50; the measured rate of 0.272 (95\% CI [0.253, 0.292]) does not satisfy this threshold, and the gate was recorded as PIVOT. A fourth sub-hypothesis (H-M3), which would have tested cross-model ranking stability using Mistral-7B-Instruct, was not executed.

\textbf{Contributions.} (1) A controlled, matched-condition comparison of token entropy, semantic entropy, and SelfCheckGPT-BERTScore on HaluEval-QA with LLaMA-2-7B-chat, to the authors' knowledge the first such comparison on this benchmark. (2) Empirical quantification of the NLI aggregation rate of microsoft/deberta-large-mnli on HaluEval-QA short factual QA responses: 0.272 (95\% CI [0.253, 0.292]), establishing a concrete domain boundary for semantic entropy applicability. (3) Documentation of a below-random AUROC for SelfCheckGPT-BERTScore (0.3562) on HaluEval-QA and a proposed, unverified explanation — label polarity inversion due to the benchmark's construction methodology. (4) A null result with mechanism analysis: each failure is decomposed into a testable mechanism claim, two of which are experimentally verified (H-M1 degenerate diagnosis, H-M2 aggregation failure).

